\begin{frame}{실습 (2/3): 비대칭의 기원 + 시드·베이스라인}
  \begin{itemize}
    \item 손계산(약 $-38$ / 약 $-20$)과 시뮬레이션($-45.3$ / $-21.3$)이
      \textbf{같은 크기·같은 순서}. Q 쪽 차이는 추락 시 깊이 분포 +
      랜덤 우회 스텝의 추가 비용.
    \item \kq{``절벽과 한 스텝'' 경로와 ``절벽과 두 스텝'' 경로는
      $\varepsilon$이 섞이는 순간 추락 확률이 \textbf{1차 vs 2차}로
      갈라서 성과 격차가 크게 벌어진다.}
  \end{itemize}
  \vskip0.5em
  \begin{enumerate}
    \setcounter{enumi}{2}
    \item \textbf{패턴 1 검증 --- 시드 2개 추가} --- 시드 1, 2로 같은
      프로토콜 재실행 $\to$ 시드별 greedy 표. (이번 실행에서는 3시드
      모두 $-13/-13/-13$, $-17/-17/-17$로 일관. 단, 8.1처럼 설정이
      조금만 달라지면 SARSA 시드 2번 $-200$ 고착도 --- ``시드 3개''가
      \textbf{형식이 아니라 정량적} 요구인 이유.)
    \item \textbf{패턴 5 검증 --- 베이스라인} --- 무작위 정책 500
      에피소드 $\to$ 평균 약 \textbf{$-5036$}(평균 486스텝, 절벽 반복
      추락). 배운 정책($-13/-17$)의 이득이 \textbf{측정된} 기준점
      위에서 확인.
    \item \textbf{리뷰 작성} --- 실습 숫자를 근거로 3점 리뷰를 템플릿으로
      완성.
  \end{enumerate}
\end{frame}
